\subsection{指数映射}

令 \(\mathrm{E}(x)\) 为 G 的一个指数映射在 0 附近的幂级数展开。令 \(\mathrm{L}(x)\) 为 G 的一个单射指数映射的逆映射在 0 附近的幂级数展开。由于任一指数映射在 0 处的切映射都是 \(\mathrm{L}(\mathrm{G})\) 的恒等映射，\(\mathrm{E}(x) \equiv x \mod \deg 2\) 且 \(\mathrm{L}(x) \equiv x \mod \deg 2\)。由于 \(\mathrm{E}(\mathrm{L}(x)) = \mathrm{L}(\mathrm{E}(x))\) 对充分接近 0 的 \(x\) 成立，形式幂级数 E 和 L 满足 \(\mathrm{E}(\mathrm{L}(\mathrm{X})) = \mathrm{L}(\mathrm{E}(\mathrm{X})) = \mathrm{X}\)。类似论证表明

\[
\mathrm{E} (t \mathrm{X}) = (\mathrm{E} (\mathrm{X})) ^ {[ t ]}, \quad \mathrm{L} (\mathrm{X} ^ {[ t ]}) = t \mathrm{L} (\mathrm{X})
\]

对 \(t \in K\)。

命题 3.

\begin{enumerate}
\def\labelenumi{(\arabic{enumi})}
\setcounter{enumi}{26}
\tightlist
\item
\end{enumerate}

\[
\mathrm{L} = \sum_ {p = 1} ^ {\infty} \frac {(- 1) ^ {p - 1}}{p} \psi_ {p}\tag{27}
\]

\[
\mathrm{E} = \sum_ {p = 1} ^ {\infty} \frac {1}{p !} \psi_ {p, p}\tag{28}
\]

（回忆 \(\psi_{p,p}\) 是 \(\psi_{p}\) 的 \(p\) 次齐次分量。）

\[
\operatorname{E} (t x) = (\operatorname{E} (x)) ^ {[ t ]}
\]

或者根据 (24)，

\[
\mathrm{E} (t x) = \sum_ {m \geqslant 0} \sum_ {1 \leqslant r \leqslant m} t ^ {r} \phi_ {r, m} (\mathrm{E} (x)).
\]

两边都是关于 \(t\) 和 \(x\) 的形式幂级数。比较 \(t\) 的一次项，得到

\[
x = \sum_ {m \geqslant 0} \phi_ {1, m} (\mathrm{E} (x)).\tag{29}
\]

现在，形式幂级数系统的逆若存在，则唯一存在（Algebra, Chapter IV, § 6, Corollary to Proposition 8）。于是

\[
\begin{array}{r l} \mathrm{L} (x) & = \sum_ {m \geq 0} \phi_ {1, m} (x) \quad \text { by   (29) } \\ & = \sum_ {p, m} \frac {(- 1) ^ {p}}{p} \psi_ {p, m} (x) \quad \text { by   (25) } \\ & = \sum_ {p} \frac {(- 1) ^ {p}}{p} \psi_ {p} (x), \end{array}
\]

从而得到 (i)。类似地，对于 \(t \neq 0\)，

\[
\begin{array}{l} \mathrm{L} (t x) = t \mathrm{L} ((t x) ^ {[ t ^ {- 1} ]}) \\ = t \mathrm{L} \left(\sum_ {m \geqslant 0} \sum_ {1 \leqslant r \leqslant m} t ^ {m - r} \phi_ {r, m} (x)\right) \quad \text { by } (24). \end{array}
\]

比较 t 的一次项，得到

\[
x = \mathrm{L} \Bigl (\sum_ {m \geqslant 0} \phi_ {m, m} (x) \Bigr)
\]

从而

\[
\begin{array}{r l} \mathrm{E} (x) & = \sum_ {m \geqslant 0} \phi_ {m, m} (x) \\ & = \sum_ {m \geqslant 0} \frac {1}{m !} \psi_ {m, m} (x) \quad \text { by   (26) }. \end{array}
\]

命题 4. G 所用的坐标图为规范坐标图的充要条件是：\(\psi_j = 0\) 对 \(j \geqslant 2\) 成立。

由命题 3，这一条件是充分的。设坐标图是规范的，并且 \(\psi_{i} = 0\) 对 \(2 \leqslant i < n\) 成立。则 \(nx = x^{[n]} = \sum_{i=0}^{n} \binom{n}{i} \psi_{i}(x) = nx + \psi_{n}(x)\)，从而 \(\psi_{n} = 0\)。因此，\(\psi_{j} = 0\) 对 \(j \geqslant 2\) 成立，这由关于 \(j\) 的归纳得到。

